\itemhead{3.2 思考 · 未编号题组（教材p.79）}{XJ-B1-C03-TH-P079-Y0868}
\label{XJ-B1-C03-TH-P079-Y0868}
{\footnotesize 湘教版教材p.79 / PDF第86页 / 3.2 思考 / 未编号题组}
\field{原题}{图 3.2-1 是湖北省武汉市某年 11 月 22 日 16 时到 23 日 16 时这 24 个小时内的气温变化图。该曲线给出了时间与气温间的函数关系。你能从中看出它的变化趋势和最高（或最低）温度吗？
\bankfig{C03-TH079}{0.78}}
\field{参考答案}{图中总体先下降后回升，期间有小幅波动；最低约 $11\,{}^\circ\mathrm C$（0时及3—4时附近），最高为次日16时约 $22.3\,{}^\circ\mathrm C$。不能把总体趋势说成全程单调。}
{\small \field{关联知识}{F04 函数三种表示法、图象与读数（教材帮p.85、86、93、94 / PDF94、95、102、103）；F05 单调性定义、证明与参数（教材帮p.97、98、100 / PDF106、107、109）；F06 最大值、最小值与值域（教材帮p.99、83、84 / PDF108、92、93）}
\field{方法与易错点}{把关键点与定义域一起标到图上，逐区间读交点和上下关系。 易错：将示意图预测当事实；把两曲线交点收进严格不等式；折线漏端点。}
\field{答案依据}{答案册未提供本栏目的独立对应答案；依据原题、教材概念补写参考回应。开放性总结不设唯一数值答案。}}
